\subsection{Descripción del dataset}\label{sec:descripcion}

La descripción completa está en el \nolinkurl{LEEME.md} del bucket, y aquí resumimos las columnas de cada tabla.

\subsubsection*{\texttt{indice\_imagenes.csv} (una fila por foto)}

\begin{table}[H]
  \centering
  \small
  \caption{Columnas de \texttt{indice\_imagenes.csv}.}
  \label{tab:col-indice}
  \begin{tabularx}{\textwidth}{>{\raggedright\arraybackslash}p{0.3\textwidth}>{\raggedright\arraybackslash}X}
    \toprule
    Columna & Descripción \\
    \midrule
    \texttt{ruta} & Ruta de la foto dentro de \nolinkurl{dataset_bcs/}. \\
    \texttt{id\_vaca} & ID de la vaca (conviene leerlo como texto). \\
    \texttt{vista} & posterior, lateral\_izq o lateral\_der. Queda vacío si la foto se descartó. \\
    \texttt{usable} & si o no. Vale no cuando la foto se descartó por estar borrosa, cortada, tapada o en una vista que no sirve. \\
    \texttt{usar\_en\_entrenamiento} & si cuando la foto es usable y su vaca tiene pliegue. \\
    \texttt{fold} & Pliegue de validación cruzada (0 a 4). Vacío si la vaca no entra. \\
    \texttt{bcs\_entrenamiento} & BCS con el que se entrena, que es el promedio de las evaluaciones de la vaca en las dos visitas. \\
    \texttt{bcs\_1}, \texttt{bcs\_2}, \texttt{bcs\_3} & Evaluaciones originales, primero las de la bitácora y después las de la hoja. \texttt{bcs\_2} y \texttt{bcs\_3} quedan vacíos si la vaca tiene menos evaluaciones. \\
    \texttt{n\_eval} & Número de evaluaciones de la vaca (1 a 3). \\
    \texttt{dif\_evaluadores} & Diferencia entre la evaluación mayor y la menor. Vacío si solo hubo una. \\
    \texttt{archivo\_jpg} & Ruta de la foto en la carpeta de trabajo intermedia. \\
    \texttt{archivo\_original} & Archivo original del que salió el JPG. \\
    \texttt{sha256} & Huella del archivo original, usada para detectar duplicados. \\
    \texttt{fecha\_hora}, \texttt{camara} & Fecha y hora, y teléfono con el que se tomó, según los metadatos. \\
    \texttt{ancho}, \texttt{alto} & Tamaño del JPG en píxeles. \\
    \texttt{profundidad} & si cuando el HEIC original traía mapa de profundidad (LiDAR). Los JPG no lo incluyen. \\
    \texttt{marcado\_por}, \texttt{marcado\_en} & Quién marcó la vista y cuándo (hora UTC). \\
    \bottomrule
  \end{tabularx}
\end{table}

Todas las fotos de una misma vaca tienen los mismos valores de BCS y de pliegue.

\subsubsection*{\texttt{vacas.csv} (una fila por vaca con fotos, 77)}

\begin{table}[H]
  \centering
  \small
  \caption{Columnas de \texttt{vacas.csv}.}
  \label{tab:col-vacas}
  \begin{tabularx}{\textwidth}{>{\raggedright\arraybackslash}p{0.3\textwidth}>{\raggedright\arraybackslash}X}
    \toprule
    Columna & Descripción \\
    \midrule
    \texttt{id\_vaca} & ID de la vaca. \\
    \texttt{en\_entrenamiento} & si cuando la vaca tiene al menos una foto usable y BCS. \\
    \texttt{fold} & Pliegue (0 a 4). \\
    \texttt{bcs\_entrenamiento}, \texttt{bcs\_1}, \texttt{bcs\_2}, \texttt{bcs\_3}, \texttt{n\_eval}, \texttt{dif\_evaluadores} & Igual que en \texttt{indice\_imagenes.csv}. \\
    \texttt{n\_fotos} & Total de fotos de la vaca, incluidas las descartadas. \\
    \texttt{n\_usables} & Fotos usables. \\
    \texttt{n\_posterior}, \texttt{n\_lateral\_izq}, \texttt{n\_lateral\_der} & Fotos usables de cada vista. \\
    \bottomrule
  \end{tabularx}
\end{table}

\subsubsection*{\texttt{evaluaciones.csv} (una fila por evaluación, 108)}

\begin{table}[H]
  \centering
  \small
  \caption{Columnas de \texttt{evaluaciones.csv}.}
  \label{tab:col-evaluaciones}
  \begin{tabularx}{\textwidth}{>{\raggedright\arraybackslash}p{0.3\textwidth}>{\raggedright\arraybackslash}X}
    \toprule
    Columna & Descripción \\
    \midrule
    \texttt{fila\_excel} & Fila de la bitácora original, para ubicar la evaluación. Vacía en las evaluaciones de la hoja. \\
    \texttt{id\_vaca} & ID de la vaca. \\
    \texttt{hora} & Hora de la evaluación (HH:MM). Vacía si la bitácora decía ``x'' o venía en otro formato, y en las evaluaciones de la hoja. \\
    \texttt{bcs} & BCS de esa evaluación, sin promediar. \\
    \texttt{orden\_evaluacion} & 1 a 3, primero en el orden de la bitácora y después en el de la hoja. \\
    \texttt{n\_eval} & Número de evaluaciones de la vaca. \\
    \texttt{fold}, \texttt{en\_entrenamiento} & Pliegue de la vaca y si entra al entrenamiento. \\
    \bottomrule
  \end{tabularx}
\end{table}

La distribución del BCS de las 108 evaluaciones es 2.0 (1), 2.25 (2), 2.50 (19), 2.75 (18), 3.00 (27), 3.25 (19), 3.50 (14) y 3.75 (8). El rango es el mismo de 2.3, de 2.0 a 3.75, así que el modelo no verá vacas fuera de ese rango, aunque el 2.0 viene de una sola vaca.
